% Part: turing-machines
% Chapter: undecidability
% Section: trakhtenbrot

\documentclass[../../../include/open-logic-section]{subfiles}

\begin{document}

\olfileid{tur}{und}{tra}
\olsection{تراختنبرُوت دا قضیہ}

\begin{explain}
  \olref[rep]{sec} وچ اسی !!{sentence}s~$!T(M,w)$
  تے~$!E(M,w)$ دی تعریف کیتی سی،
  ٹیورنگ مشین~$M$ تے انپٹ لڑی~$w$ لئی۔ فیر
  \olref[ver]{lem:valid-if-halt} تے
  \olref[ver]{lem:halt-if-valid} وچ دکھایا
  کہ~$!T(M,w) \lif !E(M,w)$ اُدوں تے صرف اُدوں
  درست اے جدوں~$M$ انپٹ~$w$ اُتے شروع
  ہو کے آخرکار رُکدی اے۔ رُکّن دا مسئلہ
  ناقابلِ فیصلہ اے۔ ایس توں نکلدا اے کہ
  پہلے درجے دی منطق وچ !!{sentence}s درست نیں یا
  نہیں، ایہ فیصلہ نہیں ہو سکدا۔ ایہ وی فیصلہ نہیں
  ہو سکدا کہ اوہ پورے ہو سکدے نیں یا نہیں
  (\cref{tur:und:uns:thm:decision-prob,tur:und:uns:cor:undecidable-sat})۔

  پر جملیاں دی درستی تے پورا ہو سکّن دیاں تعریفاں
  ہر قسم دیاں !!{structure}s لئی دتیاں
  جاندیاں نیں، چاہے اوہ محدود ہون یا لامتناہی۔
  تُہانوں شاید لگے کہ ایہ فیصلہ کرنا سوکھا
  ہووے گا کہ اک دتا ہویا !!a{sentence} کسے محدود
  !!{structure} وچ پورا ہو سکدا اے یا نہیں
  (یا ساریاں محدود !!{structure}s وچ سچ
  اے یا نہیں)۔ اسی فیصلے دے مسئلے دی
  حل نہ ہو سکّن والا ثبوت ایویں ڈھال سکدے
  آں کہ اوہ ایہناں صورتاں وچ وی فیصلے دا ناممکن ہونا دکھاوے۔

  پہلے~\olref[ver]{lem:halt-if-valid} دے
  ثبوت ول ویکھو۔ اوتھے اسی اک
  نمونہ~$\Struct M$ بنایا سی جو~$!T(M,w)$
  نوں پورا کردا اے تے جیہڑا
  عین اوہ کم بیان کردا اے جیہڑا مشین~$M$
  انپٹ~$w$ اُتے شروع ہو کے کردی اے۔
  اوس نمونے دا دائرہ~$\Nat$ سی، یعنی لامتناہی۔
  پر جے~$M$ انپٹ~$w$ اُتے واقعی رُکدی
  ہووے، تاں اسی ایس طرح اک محدود
  نمونہ~$\Struct M'$ بنا سکدے آں۔ فرض کرو
  $M$ انپٹ~$w$ اُتے شروع ہو کے~$k$ قدم
  بعد رُکدی اے۔ دائرہ~$\Domain{M'}$ لئی
  سیٹ~$\{0, \dots, n\}$ لو، جتھے
% PNB-CMP-038: leave a point strictly beyond the last executed time and needed tape position.
  $n$، $k$ تے~$w$ دی لمبائی وچوں وڈے
  عدد توں اک ودھ اے؛ تے ایہ تعبیر رکھو:
  \[
    \Assign{\prime}{M'}(x) = 
    \begin{cases}
      x + 1 &\text{جے $x < n$ ہووے}\\
      n &\text{ورنہ،}
    \end{cases}
  \]
  تے~$\tuple{x,y} \in \Assign{<}{M'}$
  اُدوں تے صرف اُدوں جدوں~$x < y$ یا~$x = y = n$۔
  باقی~$\Struct{M'}$ نوں~$\Struct{M}$ وانگ
  تعریف کرو۔ $\Struct{M'}$ دی تعریف
  مطابق، عین~\olref[ver]{lem:halt-if-valid}
  دے ثبوت وانگ، $\Sat{M'}{!T(M,w)}$۔
  تے چونکہ اسی فرض کیتا اے کہ~$M$
  انپٹ~$w$ اُتے رُکدی اے، ایس لئی
  $\Sat{M'}{!E(M,w)}$۔ لہٰذا
  $\Struct{M'}$، $!T(M,w) \land !E(M,w)$
  دا محدود نمونہ اے (دھیان دیو کہ
  اسی~$\lif$ دی تھاں~$\land$ رکھیا اے)۔
  
  ثبوت دا ادھا کم ہو گیا: اسی دکھا دتا
  کہ جے~$M$ انپٹ~$w$ اُتے رُکے، تاں
  $!T(M,w) \land !E(M,w)$ دا محدود نمونہ
  اے۔ پر ایہدی اُلٹ گل سچ نہیں۔ کجھ
  ٹیورنگ مشیناں کسے انپٹ~$w$
  اُتے نہیں رُکدیاں، پر~$!T(M,w)
  \land !E(M,w)$ دا فیر وی محدود نمونہ
  ہو سکدا اے۔ مثال لئی مشین~$M$ لو
  جیہدی اکّو حالت~$q_0$ تے اکّو
  حکم~$\delta(q_0,\TMblank) =
  \tuple{q_0,\TMblank,\TMstay}$ اے۔
  خالی انپٹ~$w = \emptyseq$
  اُتے شروع ہو کے ایہ مشین کدی نہیں
  رُکدی: ایہ لامتناہی چکر وچ اے، پر نہ
  پٹی بدلدی اے نہ سر ہلاندی اے۔
  ساریاں مکمل حالتاں اکّو جہیاں نیں
  (اوہی حالت، سر دی اوہی جگہ، پٹی دے
  اوہی نشان)۔ اسی اک محدود
  !!{structure}~$\Struct{M''}$ تعریف کر
  سکدے آں جیہڑی~$!T(M,\emptyseq) \land
  !E(M,\emptyseq)$ پورا کردی اے (مشق)۔
  البتہ اسی~$!T(M,w)$ نوں مناسب طریقے
  نال بدل سکدے آں تاکہ ایداں دیاں
  !!{structure}s خارج ہو جان۔

  \begin{prob}
    ٹیورنگ مشین~$M$ لو جیہدی
    اکّو حالت~$q_0$ تے اکّو
% PNB-CMP-039: the sole state and both successor interpretations are M-double-prime.
    حکم~$\delta(q_0,\TMblank) =
    \tuple{q_0,\TMblank,\TMstay}$ اے۔
    $\Domain{M''} = \{0, 1, 2\}$،
    $\Assign{\prime}{M''}(0) =
    \Assign{\prime}{M''}(1) = 1$
    تے~$\Assign{\prime}{M''}(2) = 2$ رکھو؛
    تے~$\Assign{<}{M''} = \{\tuple{0,1},
    \tuple{1,1}, \tuple{2,2}\}$۔
    $\Assign{\Obj Q_{q_0}}{M''}$،
    $\Assign{\Obj
    S_{\TMblank}}{M''}$ تے
    $\Assign{\Obj S_{\TMendtape}}{M''}$
    ایویں تعریف کرو کہ~$!T(M,\emptyseq)$
    تے~$!E(M,\emptyseq)$ دونیں سچ
    ہو جان، تے سمجھاؤ کیوں۔ اشارہ:
    $\delta(q_0, \TMendtape)$
    غیر تعریف شدہ اے۔ یقین کرو کہ
    \begin{align*}
    & \Obj Q_{q_0}(\num{1}, \num{n}) \land \Obj S_{\TMendtape}(\num{0}, \num{n})
      \land 
      \lforall[x][(\num{0} < x \lif \Obj S_\TMblank(x, \num{n}))] \text{\quad ہر $n \in \Nat$ لئی}\\
    & \lexists[y][(\Obj Q_{q_0}(\num{0}, y) \land \Obj S_{\TMendtape}(\num{0}, y))]
    \end{align*}
    دونیں~$\Struct{M''}$ وچ سچ ہون۔
  \end{prob}
\end{explain}

ٹیورنگ مشین~$M$ دے انپٹ
$w = \sigma_{i_1}\dots\sigma_{i_k}$ اُتے
چلن نوں بیان کرن والے !!{sentence}s ویکھو:
\begin{enumerate}
\item عدداں تے~$<$ نوں بیان کرن والے
مسلّمے (بالکل~\olref[rep]{sec} وچ
$!T(M,w)$ دی تعریف وانگ)۔
\item ابتدائی مکمل حالت نوں بیان کرن والے
مسلّمے: $!T(M,w)$ دی تعریف وانگ۔
\item اک مکمل حالت توں اگلی مکمل حالت
ول تبدیلی نوں بیان کرن والے مسلّمے:

اگے~$!A(x, y)$ پہلے وانگ ہووے، تے
\[
  !B(y) \ident \lforall[x][(x < y \lif \eq/[x][y])].
\]
\begin{enumerate}
\item \ollabel{rep-right} ہر حکم
$\delta(q_i, \sigma) =
  \tuple{q_j, \sigma', \TMright}$
لئی ایہ !!{sentence}:
\begin{align*}
& \lforall[x][\lforall[y][(
   (\Obj Q_{q_i}(x, y) \land \Obj S_{\sigma}(x, y)) \lif {}]] \\
&\qquad   (\Obj Q_{q_j}(x', y') \land \Obj S_{\sigma'}(x, y') \land
!A(x, y) \land !B(y')))
\end{align*}
\item \ollabel{rep-left} ہر حکم
$\delta(q_i, \sigma) =
  \tuple{q_j, \sigma', \TMleft}$
لئی ایہ !!{sentence}:
% PNB-CMP-040: non-edge left writes x-prime and must also force B at the next time.
\begin{align*}
& \lforall[x][\lforall[y][
    ((\Obj Q_{q_i}(x', y) \land \Obj S_{\sigma}(x', y)) \lif {}]]\\
& \qquad   (\Obj Q_{q_j}(x, y') \land \Obj S_{\sigma'}(x', y') \land
!A(x', y) \land !B(y'))) \land {}\\
& \lforall[y][((\Obj Q_{q_i}(\num{0}, y) \land \Obj S_{\sigma}(\num{0},
    y)) \lif {}]\\
& \qquad (\Obj Q_{q_j}(\num{0}, y') \land \Obj S_{\sigma'}(\num{0},
  y') \land !A(\num{0}, y) \land !B(y')))
\end{align*}
\item \ollabel{rep-stay} ہر حکم
$\delta(q_i, \sigma) =
  \tuple{q_j, \sigma', \TMstay}$
لئی ایہ !!{sentence}:
\begin{align*}
& \lforall[x][\lforall[y][(
   (\Obj Q_{q_i}(x, y) \land \Obj S_{\sigma}(x, y)) \lif {}]] \\
&\qquad (\Obj Q_{q_j}(x, y') \land \Obj S_{\sigma'}(x, y') \land
!A(x, y) \land !B(y')))
\end{align*}
\end{enumerate}
ویکھو کہ~$M$ دیاں حالتی تبدیلیاں بیان
کرن والے !!{sentence}s وچوں ہر اک
$!T(M,w)$ دے متعلقہ !!{sentence}
ورگا اے۔ ہر اک دے آخر وچ
$!B(y')$ جوڑیا اے۔ $!B(y')$ ایہ
یقینی بناندا اے کہ «اگلی» مکمل حالت
دا نمبر~$y'$، پچھلے سارے نمبراں
$\Obj 0$، $\Obj 0'$، \dots{} توں وکھرا اے۔
% \item Sentences that express consistency conditions:
% \begin{enumerate}
%   \item\ollabel{rep-con-symb}%
%   !!^a{sentence} that says that no square can have more than one
%   symbol on it at any given time:
%   \[\lforall[x][\lforall[y][(\Obj S_{\sigma}(x, y) \lif \lnot \Obj
%   S_{\sigma'}(x, y))]],\]
%   for every pair $\sigma \neq \sigma'$ of tape symbols of~$T$.
%   \item\ollabel{rep-con-state}%
%   !!^a{sentence} that says that $M$ cannot be in more than one state
%   at any given time:
%   \[\lforall[x][\lforall[y][(\Obj Q_{q}(x, y) \lif
%   \lforall[z][\lnot \Obj
%   Q_{q'}(z, y))]],\]
%   for every pair $q \neq q'$ of states of~$T$.
%   \item\ollabel{rep-con-tape}%
%   !!^a{sentence} that says that $M$ cannot be on more than one tape
%   square at any given time:
%   \[\lforall[x][\lforall[y][(\Obj Q_{q}(x, y) \lif
%   \lforall[z][\eq/[z][y]\lif \lnot \Obj
%   Q_{q'}(z, y))]],\]
%   for every pair $q$, $q'$ of states of~$T$.
% \end{enumerate}
\end{enumerate}
اُتلے سارے !!{sentence}s جوڑ کے بنے عطف نوں
$!T'(M, w)$ آکھو، ٹیورنگ
مشین~$M$ تے انپٹ~$w$ لئی۔

\begin{lem}\ollabel{lem:halts-sat}
  جے~$M$ انپٹ~$w$ اُتے شروع ہو کے
  رُکدی اے، تاں~$!T'(M,w) \land !E(M,w)$
  دا اک محدود نمونہ اے۔
\end{lem}

\begin{proof}
  $\Struct{M'}$ نوں
  \olref[ver]{lem:halt-if-valid} دے ثبوت
  وانگ لو، پر ایہ رکھو:
  \begin{align*}
    \Domain{M'} & = \{0, \dots, n\},\\
    \Assign{\prime}{M'}(x) & = 
    \begin{cases}
      x + 1 &\text{جے $x < n$ ہووے}\\
      n &\text{ورنہ،}
    \end{cases} \\
    \tuple{x,y} \in \Assign{<}{M'} &\text{عین اوہدوں جدوں $x < y$ یا $x = y = n$،}
  \end{align*}
% PNB-CMP-038: cap strictly beyond the last transition and input extent.
  جتھے~$n = \max(k,\len{w})+1$،
  تے~$k$ اوہ سب توں گھٹ عدد اے جیہدے لئی
  مشین~$M$ انپٹ~$w$ اُتے شروع ہو
  کے~$k$ قدم چُکّن توں بعد رُک
  چکی ہووے۔
  اسی ایہ جانچ مشق لئی چھڈدے آں کہ
% PNB-CMP-041: E is a named formula, and the modified theory is T-prime.
  $\Sat{M'}{!T'(M,w) \land !E(M,w)}$۔
\end{proof}

\begin{prob}
  \olref[tur][und][tra]{lem:halts-sat}
  دا ثبوت مکمل کرو: ایہ ثابت کرو کہ
% PNB-CMP-041: exercise must establish the new T-prime conjunction with !E.
  $\Sat{M'}{!T'(M,w) \land !E(M,w)}$۔
\end{prob}

\begin{lem}\ollabel{lem:sat-halts}
  جے~$!T'(M,w) \land !E(M,w)$ دا
  اک محدود نمونہ ہووے، تاں~$M$
  انپٹ~$w$ اُتے شروع ہو کے رُکدی اے۔
\end{lem}

\begin{proof}
  اسی متضادِ نقیض ثابت کراں گے۔
  فرض کرو~$M$ انپٹ~$w$ اُتے شروع ہو
  کے نہیں رُکدی۔ جے~$!T'(M,w) \land
  !E(M,w)$ دا سرے توں کوئی نمونہ
  نہ ہووے تاں کم پورا اے۔ ایس لئی
% PNB-CMP-042: the assumed model belongs to the lemma's T-prime theory.
  فرض کرو~$\Struct{M}$،
  $!T'(M,w) \land !E(M, w)$ دا نمونہ
  اے۔ سانوں دکھانا اے کہ ایہ محدود
  نہیں ہو سکدا۔
  
  اسی~\olref[ver]{lem:config} وانگ
  ثابت کر سکدے آں کہ جے~$M$،
  انپٹ~$w$ اُتے شروع ہو کے
  $n$ قدم بعد ہن تک نہ رُکی ہووے
  تے ایہ قدم-گنتی مثبت ہووے،
% PNB-CMP-043: B is forced at positive successor times, not the zero base.
  تاں~$!T'(M,w)
  \Entails !C(M, w, n) \land
  !B(\num{n})$؛ صفر اُتے ابتدائی
  مکمل حالت سِدھی ابتدائی مسلّماں
  توں ملدی اے۔ چونکہ~$M$ انپٹ~$w$
  اُتے شروع ہو کے کدی نہیں رُکدی،
  ایس لئی~$!T'(M,w) \Entails !C(M, w, n) \land
  !B(\num{n})$ ہر مثبت~$n \in \Nat$
  لئی۔ دھیان دیو کہ
  \olref[rep]{prop:mlessk} مطابق
  $!T'(M,w) \Entails \num{k} < \num{n}$
  ہر~$k < n$ لئی۔ نالے
  $!B(\num{n}) \Entails \num{k} < \num{n} \lif
  \eq/[\num{k}][\num{n}]$۔
  ایس لئی~$\Sat{M}{\eq/[\num{k}][\num{n}]}$
  ہر~$k < n$ لئی،
  یعنی ٹرماں~$\num{k}$ دی لامتناہی
  لڑی دیاں قدراں~$\Struct{M}$ وچ
  آپس وچ وکھ وکھ ہون گیاں۔ ایس لئی
  $\Domain{M}$ لامتناہی ہونا
  ضروری اے؛ ایس لئی~$\Struct{M}$،
  $!T'(M,w) \land !E(M, w)$
  دا محدود نمونہ نہیں ہو سکدا۔
\end{proof}

\begin{prob}
  \olref[tur][und][tra]{lem:sat-halts}
  دا ثبوت مکمل کرو: ایہ ثابت کرو کہ
  جے~$M$ انپٹ~$w$ اُتے شروع ہو کے
  $n$ قدم بعد نہ رُکی ہووے، تاں
% PNB-CMP-043: the exercise needs the positive-time qualification.
  $!T'(M,w) \Entails !B(\num{n})$
  ہر مثبت عدد لئی۔ صفر دی
  مکمل حالت ابتدائی مسلّماں توں لو۔
\end{prob}

\begin{thm}[تراختنبرُوت دا قضیہ]
  \ollabel{thm:trakhtenbrodt}
  ایہ فیصلہ نہیں کیتا جا سکدا کہ
  پہلے درجے دی منطق دا دتا ہویا
  !!{sentence} کوئی محدود نمونہ رکھدا اے
  یا نہیں (یعنی اوہ کسے محدود
  ساخت وچ پورا ہو سکدا اے یا نہیں)۔
\end{thm}

\begin{proof}
  فرض کرو کوئی ٹیورنگ مشین~$F$
  محدود نمونہ ہون دا فیصلہ
  کردی ہووے۔ فیر کسے وی
  ٹیورنگ مشین~$M$ تے انپٹ~$w$
  توں اسی~$!T'(M,w) \land !E(M,w)$
  جملہ حساب کرکے~$F$ نال فیصلہ
  کر سکدے آں کہ ایہدا محدود
  نمونہ اے یا نہیں۔
  \cref{tur:und:tra:lem:halts-sat,tur:und:tra:lem:sat-halts}
  مطابق ایہدا محدود نمونہ اُدوں تے
  صرف اُدوں اے جدوں~$M$ انپٹ~$w$
  اُتے شروع ہو کے رُکدی اے۔
  ایس لئی~$F$ نال رُکّن دا
  مسئلہ حل ہو جائے گا۔ سانوں
  پتا اے کہ ایہ حل نہیں ہو سکدا۔
\end{proof}

\begin{cor}\ollabel{cor:fproof-incomp}%
% PNB-CMP-044: the contradiction uses effective enumeration of derivations.
  ہر \emph{محدود} ساخت وچ درستی لئی
  کوئی صحت مند تے کامل !!{derivation} دا نظام،
  جیہدے اخذ مؤثر طور اُتے گنتی وار
  فہرست وچ لکھے جا سکن، نہیں ہو سکدا۔
  ایہ ایہو جیہا !!a{derivation} دا نظام ہووے گا
  جیہدے وچ~$\Proves !B$
  اُدوں تے صرف اُدوں ہووے جدوں
  $\Sat{M}{!B}$ ہر محدود
  !!{structure}~$\Struct{M}$ وچ ہووے۔
\end{cor}

\begin{proof}
  مشق۔
\end{proof}

\begin{prob}
  \olref[tur][und][tra]{cor:fproof-incomp}
  ثابت کرو۔ دھیان دیو کہ
  $!B$ ہر محدود !!{structure}
  وچ اُدوں تے صرف اُدوں پورا اتردا
  اے جدوں~$\lnot !B$ کسے وی محدود
  ساخت وچ پورا نہیں ہو سکدا۔
  سمجھاؤ کہ محدود نمونہ ہون دا مسئلہ،
  \olref[tur][und][uns]{thm:valid-ce}
  والے معنی وچ نیم فیصلہ
  کرن جوگ کیوں اے۔ ایس
  توں دلیل دیو کہ جے محدود
% PNB-CMP-044: effective enumeration is needed for the corollary's decision procedure.
  درستی لئی کوئی مؤثر
  !!a{derivation} دا نظام
  ہووے، تاں محدود نمونہ ہون دا
  فیصلہ ہو سکدا ہووے گا۔
\end{prob}

\end{document}
